\documentclass[11pt,a4paper]{article}

%%%%%%%%%%%%%%%%%%%%%%%%%%%%%%%%%%%%%%%%%%%%%%%%%%%%%%%%%%%%%%%%%%%%%%%%%%%%
%% Standalone revision plan — data-dependent MUST items (JPGN Major Revision)
%% Self-contained: compile with either pdflatex or xelatex, e.g.
%%     pdflatex revision_plan_data_dependent_musts.tex
%% No bibliography, no external files required.
%%%%%%%%%%%%%%%%%%%%%%%%%%%%%%%%%%%%%%%%%%%%%%%%%%%%%%%%%%%%%%%%%%%%%%%%%%%%

\usepackage[margin=1in]{geometry}
\usepackage{amsmath}
\usepackage{booktabs}
\usepackage{longtable}
\usepackage{enumitem}
\usepackage{xcolor}
\usepackage{framed}
\usepackage{titlesec}
\usepackage[T1]{fontenc}
\usepackage{lmodern}
\usepackage{microtype}

% Colored links (blue, no boxes) — hyperref loaded late
\definecolor{citeblue}{RGB}{0,0,255}
\usepackage[colorlinks=true, allcolors=citeblue, breaklinks=true]{hyperref}

% ---- Light styling for the two recurring block types (framed pkg only) ----
\definecolor{paperbg}{RGB}{244,244,240}
\definecolor{refbg}{RGB}{238,242,248}
\definecolor{rulegray}{RGB}{170,170,170}

% A shaded environment with a thin left rule, colored by \FrameColor/\ShadeColor.
% "What the paper says" block — ivory background
\newenvironment{papersays}%
  {\def\FrameCommand{{\color{rulegray}\vrule width 1.5pt}\hspace{6pt}%
     \colorbox{paperbg}}%
   \MakeFramed{\advance\hsize-\width \FrameRestore}\small}%
  {\endMakeFramed}

% "Referee says" block — light blue background
\newenvironment{refereesays}%
  {\def\FrameCommand{{\color{rulegray}\vrule width 1.5pt}\hspace{6pt}%
     \colorbox{refbg}}%
   \MakeFramed{\advance\hsize-\width \FrameRestore}\small\itshape}%
  {\endMakeFramed}

\titleformat{\section}{\large\bfseries}{\thesection.}{0.5em}{}
\setlength{\parindent}{0pt}
\setlength{\parskip}{4pt}

\title{\textbf{Revision Plan --- Data-Dependent MUST Items}\\[2pt]
\large JPGN Major Revision}
\author{Chronic Intestinal Failure in Brazil: A First Estimate of Prevalence}
\date{2026-07-13}

\begin{document}
\maketitle

\noindent
This document lists the MUST items from the JPGN Major Revision that \textbf{cannot be
done from the manuscript alone}: each requires re-running or extending the analysis on the
underlying SIH-SUS data. The \texttt{scripts/} directory is currently empty, so the analysis
that produced 3.35/million, the 971-patient cohort, and all figures is not yet in the repo.
Every item below is blocked on bringing that data/code into the project.

The text-only fixes (title, linkage-precedent softening, hidden-mortality down-tone,
race/equity reframing, ``alta melhorada'' definition, appendix placeholder) were already
applied in commit \texttt{ed3d7e0} and are not repeated here.

\medskip
\hrule
\medskip

\newcommand{\papersaysHead}{\textit{What the paper currently says.}}
\newcommand{\refereeHead}{\textit{What the referees said.}}
\newcommand{\whyHead}{\textbf{Why data is required.}}
\newcommand{\opHead}{\textbf{Data operation.}}

%======================================================================
\section{Reconcile the 1,028-vs-971 count; define the estimand}
%======================================================================

\papersaysHead\ Results (line 153):
\begin{papersays}
``these admissions corresponded to \textbf{150,078 unique patients} \ldots A total of
\textbf{971 patients} (0.65\%) had at least one admission with 60 or more daily units of PN,
meeting the operational definition of CIF''
\end{papersays}
But Table~2 (\texttt{tab:prevalencia\_17anos}) lists annual ``No.\ of SUS cases'' as
148, 136, 147, 112, 86, 106, 108, 105, 80 --- which \textbf{sum to 1,028} [verified], while the
text's unique-patient total is \textbf{971}.

\refereeHead\ Methods (m1):
\begin{refereesays}
``the nine annual SUS-case counts in Table~2 \ldots sum to \textbf{1,028}, but the text
reports \textbf{971 unique CIF patients} \ldots This is defensible but must be stated
explicitly --- otherwise the table appears arithmetically inconsistent with the headline N.''
\end{refereesays}
Domain (M3b):
\begin{refereesays}
``a \textit{period prevalence} built by summing annual counts double-counts persistent
patients across years, which is conceptually inconsistent with `unique patients.'\,''
\end{refereesays}

\whyHead\ The 57-patient gap is presumably patients with qualifying admissions in
$\geq 2$ calendar years (counted once in 971, once per year in Table~2), but this cannot be
confirmed from the manuscript.

\opHead
\begin{enumerate}[leftmargin=*, topsep=2pt]
  \item Query the patient-level dataset: how many unique patients have qualifying admissions
        in $\geq 2$ calendar years? Confirm it explains $1{,}028 - 971 = 57$.
  \item Decide and state the \textbf{estimand} (period vs point prevalence); add one sentence
        to Methods $+$ a Table~2 footnote.
\end{enumerate}

%======================================================================
\section{Add confidence intervals everywhere}
%======================================================================

\papersaysHead\ Every figure is a bare point estimate. Results ($\sim$line 155):
\begin{papersays}
``prevalence fell to 3.26 and 2.53 cases per million, reductions of approximately
\textbf{23\% and 40\%} relative to the pre-pandemic mean''
\end{papersays}
Table~2 reports ``Prevalence (per million)'' as single numbers (4.20 \ldots 2.46), no interval.

\refereeHead\ Methods (M1, ``the single most serious defect''):
\begin{refereesays}
``a Poisson 95\% CI on 80 observed cases is roughly \textbf{$\pm$22\%}. A clinical readership
will not accept year-over-year `declines' asserted from counts this small without at least
Poisson/exact intervals.''
\end{refereesays}
Domain (M2, ``non-negotiable for this journal''):
\begin{refereesays}
``Report exact (Poisson/binomial) confidence intervals for each annual prevalence and for
mortality.''
\end{refereesays}

\whyHead\ CIs are computed quantities. Annual prevalence CIs need the case count $+$
denominator (both in Table~2); the pandemic ``23\%/40\% decline'' should become a
\textbf{rate ratio with CI}; mortality CIs need the CIF-cohort numerator/denominator (see \S5).

\opHead\ R: \texttt{poisson.test} / \texttt{epitools::pois.exact} for each annual prevalence
and mortality; rate-ratio-with-CI for the pandemic comparison. Regenerate Table~2 with
interval columns.

%======================================================================
\section{Rebuild the headline as a scaling sensitivity band (editor's \#1 priority)}
%======================================================================

\papersaysHead\ Methods, scaling equation:
\begin{papersays}
``the total estimated number of cases in year $t$ was obtained by scaling the observed SUS
count by a factor of \textbf{1.55}:\quad $N_{\text{total},t} = N_{\text{SUS},t} \times 1.55$''
\end{papersays}
This single scalar produces the headline \textbf{3.35/million}.

\refereeHead\ This is the \textbf{FATAL-as-presented} flag. Methods (M2):
\begin{refereesays}
``the entire non-SUS column of Table~2 is mechanically $0.55 \times N_{\text{SUS}}$ \ldots the
reader cannot separate `signal' from `this one scalar' \ldots 1.55 may \textit{overstate} the
private share --- meaning the direction of bias is not even signed.''
\end{refereesays}
Editor (single most important revision):
\begin{refereesays}
``Reconstruct the headline as an \textit{honest interval rather than a point estimate} \ldots
anchor it on the unscaled SUS-only floor, flank the 1.55 with a sensitivity band.''
\end{refereesays}

\whyHead\ Prevalence must be recomputed under a grid of scaling assumptions.

\opHead
\begin{itemize}[leftmargin=*, topsep=2pt]
  \item $\times 1.0$ (unscaled SUS-only) $=$ strict lower bound $\rightarrow$ promote to
        \textit{primary} estimate.
  \item $\times 1.55$ $=$ current central estimate.
  \item Upper bound from the all-cause PNS split at its plausible range.
  \item Ideally anchor a \textit{disease-specific} SUS share (where RESTORE/Leite/Goldani
        Brazilian cohorts were treated --- predominantly public centers).
  \item Report the national estimate as a \textbf{range}; regenerate Table~2 / the headline.
\end{itemize}

%======================================================================
\section{Validate that ``di\'{a}rias'' (daily units) $=$ days; resolve the 288 contradiction}
%======================================================================

\papersaysHead\ Definition (Methods):
\begin{papersays}
``parenteral nutrition (PN) support for at least \textbf{60 days} within a consecutive
74-day interval''
\end{papersays}
Operationalization:
\begin{papersays}
``we retained only those in which the recorded duration of the PN procedure was
\textbf{60 daily units (di\'{a}rias) or more}''
\end{papersays}
Table~1 reports a \textbf{``Max daily units'' of 288 in 2020} --- which exceeds a 74-day window.

\refereeHead\ Domain (M1, ``the single most important validity question for the whole paper''):
\begin{refereesays}
``In SIH-SUS the `di\'{a}ria' is a billing/reimbursement unit, not a guaranteed
one-per-calendar-day count. If a patient can accrue more than one PN di\'{a}ria per day \ldots
then the count is not a clean day-count and the 60-di\'{a}ria threshold does not map 1:1 to the
60-day ASPEN criterion.''
\end{refereesays}
Methods (m5):
\begin{refereesays}
``Table~1 `Max daily units' of 288 in 2020 \ldots \textbf{exceeds the number of days that could
plausibly accrue in a 74-day window}; if daily units accumulate across a full year or across
linked admissions, the 74-day-window definition is again ambiguous.''
\end{refereesays}

\whyHead\ Requires the DATASUS codebook $+$ inspection of raw records.

\opHead
\begin{enumerate}[leftmargin=*, topsep=2pt]
  \item Confirm from the DATASUS data dictionary whether one AIH ever bills $>1$ PN
        di\'{a}ria/calendar day; cite it (or correct the mapping). The 288 value suggests
        di\'{a}rias accumulate across multiple linked admissions.
  \item State precisely, from the actual code, whether the 60/74 criterion is applied
        \textbf{within a single admission} or \textbf{cumulatively across linked admissions},
        and how the 74-day window is enforced.
\end{enumerate}

%======================================================================
\section{Fix the mortality inconsistency (which figure is the CIF cohort's?)}
%======================================================================

\papersaysHead\ Three different numbers a reader takes as the same thing:
\begin{papersays}
\begin{itemize}[leftmargin=1.2em, topsep=1pt, itemsep=1pt]
  \item \textbf{Abstract:} ``In-hospital mortality ranged from \textbf{9\% to 15\%}''
  \item \textbf{Results, 971-cohort (line 153):} ``\textbf{6.5\%} died during hospitalization''
  \item \textbf{Table~1} (\texttt{tab:pn\_17anos}, ``\% Deaths''): 11.9\%--15.0\% --- but for
        \textbf{all 198,769 PN admissions}, not the 971 CIF patients.
\end{itemize}
\end{papersays}

\refereeHead\ Domain (M3a):
\begin{refereesays}
``the `9--15\%' quoted for CIF appears to be borrowing the all-PN mortality column \ldots a
reader cannot tell whether CIF-cohort mortality is $\sim$6.5\% or $\sim$9--15\%, and the
Discussion's comparison to high-income survival ($\geq$85\%) hinges on getting this right.''
\end{refereesays}

\whyHead\ The correct cohort-specific mortality must be computed from the patient-level data.
Figure \texttt{fig:obito} appears to show all-PN mortality, not CIF-cohort mortality, and may
need regenerating.

\opHead\ Compute in-hospital mortality for the 971-patient CIF cohort (overall $+$ annual for
\texttt{fig:obito}); use that single figure consistently in abstract, Results, Discussion.

%======================================================================
\section{Quantify the missing-CEP exclusion and test for regional selectivity}
%======================================================================

\papersaysHead\ Methods:
\begin{papersays}
``Admissions with missing postal code information were excluded from the analysis, as the
absence of this field precludes reliable patient matching.''
\end{papersays}
No count is given for how many admissions/patients this removes.

\refereeHead\ Methods (M5, ``potentially large''; the most consequential distinct concern):
\begin{refereesays}
``missing CEP is unlikely to be random --- it plausibly correlates with the
North/Northeast/Central-West underserved regions the paper's central argument concerns. If
missingness is concentrated in exactly the regions claimed to be underrepresented,
\textbf{the exclusion mechanically reinforces the paper's headline conclusion}, and the
`lower bound' framing may be partly an artifact of the exclusion.''
\end{refereesays}

\whyHead\ Purely empirical; requires the pre-exclusion dataset.

\opHead
\begin{enumerate}[leftmargin=*, topsep=2pt]
  \item Count admissions/patients dropped for missing CEP --- \textbf{overall and by region}.
  \item Compare CEP-present vs CEP-absent admissions on observables (region, age, death).
  \item A best/worst-case bound on the headline under assumptions about the excluded cases.
\end{enumerate}

%======================================================================
\section{(Data half) False-match / twin-collision rate estimate}
%======================================================================

\textit{Note.} The \textit{text} half (softening the linkage-precedent overclaim, flagging twin
risk) was already applied in commit \texttt{ed3d7e0}. The \textbf{data half} remains.

\refereeHead\ Methods (M4):
\begin{refereesays}
``Twins are the specific concern: same CEP, same DOB, often same sex --- the deterministic rule
will merge them into one `patient,' \textbf{biasing the count downward} in exactly the
prematurity-driven population that constitutes the sample.''
\end{refereesays}
Suggested internal validation:
\begin{refereesays}
``agreement between the deterministic linkage and the CIF-specific-code deduplication on the
22 code-identified patients''
\end{refereesays}

\whyHead\ Both are direct data cross-tabulations.

\opHead
\begin{enumerate}[leftmargin=*, topsep=2pt]
  \item Count CIF triplets (CEP $+$ sex $+$ DOB) that plausibly represent multiple births
        merged into one record (same-DOB/CEP/sex clusters); bound the count treating each as
        1 vs 2 patients.
  \item \textbf{22-patient overlap check (SHOULD-1):} of the 22 CIF-code patients, how many
        does the $\geq$60 PN rule independently flag?
\end{enumerate}

%======================================================================
\section*{Summary table}
%======================================================================
\addcontentsline{toc}{section}{Summary table}

\begin{longtable}{@{}c p{4.3cm} p{6.6cm}@{}}
\toprule
\textbf{MUST} & \textbf{One-line need} & \textbf{Data operation} \\
\midrule
\endhead
1 & Explain 1,028 vs 971 & Count patients appearing in $\geq$2 years \\
2 & Confidence intervals & Poisson/binomial CIs on counts ($+$ rate ratio for pandemic) \\
3 & Scaling band & Recompute prevalence at $\times$1.0 / $\times$1.55 / upper \\
4 & di\'{a}rias $=$ days? & Inspect codebook $+$ confirm windowing logic \\
5 & Correct mortality & Compute mortality on the 971 cohort \\
6 & Missing-CEP bias & Count exclusions by region; compare observables \\
7 (data) & Twin/false-match rate & Count same-DOB/CEP/sex clusters; 22-patient overlap \\
\bottomrule
\end{longtable}


\end{document}
